% 原書p.364 constraintからp.369 first-quantized descriptionの項末まで。
\gloss{拘束（constraint）}
{BRSTまたはOCQの意味で、物理的状態における行列要素がゼロであることを要求される対称性の生成子。通常はゲージ対称性から生じると理解できる。これは、いわゆる第1種拘束についてのことであり、本書で考えるのはそれだけである。拘束は、ゲージ対称性の生成子のうち、運動方程式によっては消えないものからなる。}

\gloss{剰余CFT（coset CFT）}
{既知のCFTのエネルギー運動量テンソルが、互いに可換な部分の和として書けるとき、そのCFTの因子の1つとして構成されるCFT。古典的な例では、完全なCFT $G$と因子の1つ$H$が、ともにカレント代数である。これは対称性$H$をゲージ化することとも考えられる。}

\gloss{宇宙定数（cosmological constant）}
{真空のエネルギー密度。超対称でない量子論には、超対称性が自発的に破れたものも含めて、このようなエネルギー密度を生じる効果が数多くある。\emph{宇宙定数問題}とは、自然界の宇宙定数が、既知のゼロでない効果よりも何桁も小さいという問題である。}

\gloss{臨界次元（critical dimension）}
{平坦なMinkowski時空において摂動的な弦理論が整合的となる次元。ボソン弦では26、超対称な弦理論では10である。}

\gloss{カレント代数（current algebra）}
{量子場理論において、連続的な対称群$g$に伴うカレント、またはそのFourierモードの代数。本書で用いる場合には、エネルギー運動量テンソルをSugawara形式と定義した、2次元CFTに現れる特定の代数を指す。この代数には、アフィンLie代数、アフィンKac--Moody代数という用語も用いられる。ただし、カレント代数と同様に、どちらの用語にもより広い定義がある。アフィンという語は、$c$数のSchwinger項を指す。\emph{ねじれのないカレント代数}は、整数のモードを持つ周期的なカレントの代数である。代数は、$g$の任意の自己同型によって\emph{ねじる}ことができる。}

\gloss{サイクル（cycle）}
{ホモロジーの意味で、位相的に非自明な部分多様体。$p$サイクルは$p$次元である。$A$サイクルと$B$サイクルは、Riemann面上の非自明な1サイクルの標準的な基底である。}

\gloss{Dブレーン（D-brane）}
{I型、IIA型、IIB型の弦理論において、弦が終端できる力学的な物体。名称はDirichlet braneの短縮形である。取り付いた弦の座標は、ブレーンに垂直な方向ではDirichlet境界条件を、接する方向ではNeumann条件を満たす。D$p$ブレーンは$p$次元であり、$p$はIIA理論では任意の偶数、IIB理論では任意の奇数、I型理論では1、5、9を取る。D9ブレーンは空間を満たすので、通常のNeumann境界条件に対応する。D$p$ブレーンは$(p+1)$形式のR--Rゲージ場の源である。Dブレーンの質量または張力は、通常の量子や基本弦のものと、ソリトンのものの中間にある。Dブレーンの低エネルギーの揺らぎは超対称ゲージ理論で記述され、ブレーンが重なると、この理論は非Abel型になる。}

\gloss{Dインスタントン（D-instanton）}
{空間だけでなくEuclid時間にも局在する物体。取り付いた弦のすべての座標にDirichlet条件を課すことで定義される。場の理論のインスタントンに似ており、R--R場の強さの値を変えるトンネル過程に対応する。より一般には、D$p$ブレーンの$(p+1)$次元の世界体積が時間方向に局在し、空間内の$(p+1)$サイクルに巻き付いているときも、同様の効果を持つ。}

\gloss{D弦（D-string）}{I型およびIIB型弦理論におけるD1ブレーン。}

\gloss{$D$項（$D$-term）}
{1. 4個または8個の超対称性を持つゲージ理論における、ゲージ多重項の補助場。2. この補助場の2乗に比例するポテンシャル項。ゲージ結合だけに依存し、$U(1)$の場合にはFayet--Iliopoulosパラメーターの値にも依存する。}

\gloss{Del Guidice--Di Vecchia--Fubini（DDF）演算子}
{振動子代数を満たし、OCQで物理的状態の完全な組を生成する演算子。}

\gloss{子孫状態（descendant）}
{最高ウェイト状態にVirasoroの上昇生成子を作用させて得られる状態。}

\gloss{対角モジュラー不変量（diagonal modular invariant）}
{左向きと右向きの場に共通の境界条件を課すことで作られる、モジュラー不変なCFT。}

\gloss{$\mathrm{diff}$不変性（diff invariance）}
{一般座標不変性、すなわち再パラメーター化不変性。通常は世界面の座標に適用される。}

\gloss{ディラトン（dilaton）}
{すべての摂動的な弦理論に見いだされる、重力と同程度の強さの結合を持つ質量のないスカラー。厳密に質量のないディラトンは、重力以外の相互作用に対する制限に抵触する。しかし、ディラトンの質量を禁じる対称性はないので、超対称性が破れた真空では力学的効果によって質量が生じる。他のモジュライについても同じである。弦の結合定数はディラトン場の値によって決まる。}

\gloss{次元還元（dimensional reduction）}
{最も単純な場合には、コンパクト方向の運動量がゼロの状態だけを残すトーラス・コンパクト化。より一般には、そして物理的にはやや不明瞭な意味では、すべての場にある一組の対称性のもとでの不変性を要求することによって、低次元の場の理論を構成すること。この場合、コンパクト化による解釈が存在するとは限らない。}

\gloss{Diracの量子化条件（Dirac quantization condition）}
{電荷と磁荷について、その積が量子化されるという条件$\mu_e\mu'_m=2\pi n$。両方の電荷を持つ2つのダイオンについては、条件は\mbox{$\mu_e\mu'_m-\mu_m\mu'_e=2\pi n$}となる。これらの条件は、$d$次元における、次元が$p$および$d-p-4$の物体へ一般化される。一方は$(p+2)$形式の場の強さの源であり、他方はPoincaré双対な$(d-p+2)$形式の場の強さの源である。}

\gloss{Diracスピノル（Dirac spinor）}
{Dirac行列の代数、別名Clifford代数の一意的な既約表現。これはLorentz群の表現でもある。偶数次元では、Lorentz群の2つのWeyl表現に可約である。Diracスピノルは複素数であり、特定の次元では、さらにMajorana条件、すなわち実条件を課すことがLorentz不変性と両立する。}

\gloss{Dirichlet境界条件（Dirichlet boundary condition）}
{境界で場の値を固定する条件。弦理論で関係するのはこの用法であるが、別の文脈では接線方向の微分だけを固定すればよい場合もある。}

\gloss{離散的捩率（discrete torsion）}
{ねじれたCFTを構成する際に、経路積分の各セクターの位相を変え、それによってHilbert空間への射影を変えること。}

\gloss{二重化の手法（doubling trick）}
{境界を持つ多様体上の正則場と反正則場を、境界で反射して得られる多様体の2重のコピーの上で、正則場だけによって表すこと。}

\gloss{双対共鳴模型（dual resonance model）}
{弦理論へと発展した、強い相互作用の現象論的模型。ここで双対という語は、世界面双対性を指す。}

\gloss{双対性（duality）}
{一見異なる物理系の同値性。このような同値性は、単一の量子論が異なる古典極限を持つときにしばしば現れる。古典的な例の1つは粒子・波動双対性であり、量子場理論は、古典的な場の理論で記述される極限と、古典的な粒子力学で記述される別の極限を持つ。もう1つはIsing模型の高温・低温双対性である。ここでは、低温は古典極限の統計力学的な類似物であり、Boltzmann和は最低エネルギーの配位によって支配される。Montonen--Olive双対性、$S$双対性、弦・弦双対性、$T$双対性、$U$双対性、世界面双対性を参照。}

\gloss{有効場の理論（effective field theory）}
{与えられたエネルギースケールより低い領域、同値な言い方をすれば、与えられた長さスケールより長い領域における物理系の記述。}

\gloss{Einstein計量（Einstein metric）}
{低エネルギーでの主要な作用が、他の場に依存しないHilbert作用
\[
\frac{1}{2\kappa^2}\int d^d x\,(-G)^{1/2}\bm R
\]
となる計量。ここで$\kappa$は重力結合であり、Planck長と$\kappa=(8\pi)^{1/2}L_{\rm P}$によって関係づけられる。この計量は、場に依存するWeyl変換によって、シグマ模型の計量などの他の計量に関係づけられる。異なる計量の存在は等価原理を破るように思われるが、ディラトンや他のモジュライが質量を持つときには、その区別は消える。}

\gloss{電弱スケール（electroweak scale）}
{電弱対称性の破れの質量スケール。およそ$10^2\ \mathrm{GeV}$。}

\gloss{増大したゲージ対称性（enhanced gauge symmetry）}
{モジュライ空間の特別な点に現れる、理論のもとの定式化からは明らかでないゲージ対称性。古典的な例は、トーラス・コンパクト化の特別な半径に現れるゲージ対称性であり、そのゲージボソンは巻き付き状態である。現在では、ほかにも多くの機構が知られている。収縮するサイクルに巻き付いたDブレーンやブラック$p$ブレーン、さまざまなブレーンの間に張られたF弦やD弦でブレーン同士が重なる極限、そして大きさゼロの$SO(32)$インスタントン上に現れるゲージ対称性である。}

\gloss{Euclid的（Euclidean）}
{計量が正定値であること。計量が平坦であるというもとの意味合いは、ある程度無視されている。そのため、Minkowski的な理論の解析接続として予想される理論をEuclid量子重力と呼ぶ。計量そのものについては、Euclid符号を持つ曲がった計量に対してRiemann的という語を用いるほうが正確である。}

\gloss{Euclid随伴（Euclidean adjoint）}
{Euclid量子論において、Hermite随伴と時間反転を組み合わせたもの。後者の操作は随伴の時間反転の効果を打ち消すため、組み合わせた操作は局所的になる。}

\gloss{Euler数（Euler number）}
{位相不変量
\[
\chi=\sum_{p=0}^{d}(-1)^pB_p,
\]
ここで$B_p$は第$p$ Betti数である。種数$g$のRiemann面では$2(1-g)$に等しい。より正確には、Euler標数という。}

\gloss{期待値（expectation value）}
{指定された挿入を持つ経路積分。本書で採用した用語であるが、相関関数（correlation function、correlator）もよく使われる。}

\gloss{拡張カイラル代数（extended chiral algebra）}
{CFTの正則演算子すべての集合。}

\gloss{拡張超対称性（extended supersymmetry）}
{与えられた時空次元の最小スピノル表現を、超電荷が複数組含む超対称代数。}

\gloss{$F$項（$F$-term）}
{1. $d=4,\ N=1$超対称性のカイラル多重項の補助場。2. この場の2乗に比例するポテンシャル項。}

\gloss{F理論（F theory）}
{1. ディラトンとR--Rスカラーの非自明な背景を持つIIB超弦の状態を、これらの場を補助的な2次元トーラスのモジュラス$\tau$に関係づけることによって体系的に記述するもの。2. IIB弦の基礎にあると予想される12次元の量子論。この名称はM理論に着想を得たものであり、Fはfatherを表す。}

\gloss{Faddeev--Popov行列式（Faddeev--Popov determinant）}
{ゲージ不変な汎関数積分をゲージ断面上の積分に還元するときに生じるJacobi行列式。}

\gloss{Faddeev--Popovゴースト（Faddeev--Popov ghosts）}
{Faddeev--Popov行列式を汎関数積分で表すために使われる、通常と逆の統計に従う量子場。}

\gloss{Fayet--Iliopoulos項（Fayet--Iliopoulos term）}
{4個または8個の超対称性を持つ$U(1)$ゲージ理論において、補助場である$D$項の場について線形な作用の項。}

\gloss{Feigin--Fuchs表現（Feigin--Fuchs representation）}
{極小模型の期待値を自由場によって表す表現。}

\gloss{ファイブレーション（fibration）}
{局所的にはファイバー$F$と底空間$B$の積である空間。底空間を移動するとファイバーの幾何は変化し、特異になることもある。典型的なファイバーは、トーラスとK3多様体である。}

\gloss{第1 Chern類（first Chern class）}
{複素多様体上のRicci形式$R_{i\bar\jmath}\,dz^i\,d\bar z^{\bar\jmath}$のDolbeaultコホモロジー類。}

\gloss{第一量子化による記述（first-quantized description）}
{量子化された粒子理論を粒子の経路についての和として、または弦理論を世界面についての和として表すこと。\emph{第二量子化}は、通常の場または弦の場についての汎関数積分による表現を指す。第二量子化という用語は、第一量子化の波動関数を場の演算子として解釈し直すことを含意する。この呼び方は広く使われているが、存在しないかもしれない深い原理を示唆する点で不適切だという議論もある。世界面についての和自体が量子場理論なので、それを第二量子化と呼ぶことも同じようにでき、その場合、弦の場の理論は第三量子化となる。通常の場の理論の第三量子化は、宇宙を生成・消滅する演算子を記述することになるが、これは役に立つかもしれないし、そうでないかもしれない概念である。}
